% PDF-only preamble (pandoc -H). Reader feedback: at 11pt, mono vs serif
% alone is hard to tell apart for older eyes. Set inline code (and plain
% verbatim blocks) in a dark blue -- the PicoMite-manual convention -- so
% commands stand out from body text by colour, not just letterform.
% Syntax-highlighted (```python) blocks keep their tango colours.
%
% The colour is given numerically, not as a named colour: inline code in
% section titles reaches the running page headers, which the book class
% passes through \MakeUppercase -- a colour NAME gets uppercased into an
% undefined one, but digits survive.
\usepackage{xcolor}
\let\pcTexttt\texttt
\renewcommand{\texttt}[1]{\textcolor[RGB]{0,70,160}{\pcTexttt{#1}}}
\makeatletter
\def\verbatim@font{\ttfamily\color[RGB]{0,70,160}}
\makeatother

% Reader feedback: code listings pasted from the PDF arrived flush-left --
% their indentation was gone. Cause: xelatex stores verbatim spaces as
% horizontal positioning, not as characters, so copy/paste (and the REPL)
% never see them. Fix: make fancyvrb emit a real space glyph (\char"20) for
% every code-block space. It has the correct monospace width, so the page
% looks identical, but the spaces are now extractable and listings paste with
% their indentation intact. (\XeTeXgenerateactualtext alone does NOT fix this
% -- the spaces are glue between runs, which ActualText spans don't capture.)
\makeatletter
\def\FV@Space{\char"20\relax}
\makeatother
